\documentclass[12pt]{exam}
\pagestyle{headandfoot}
\firstpageheader{}{}{}
\runningheader{\date}{}{\class, \teacher}
\runningheadrule
\firstpagefooter{}{}{}
\runningfooter{}{Viel Erfolg!}{Seite \thepage\:von \numpages}
\runningfootrule
%=========================================================
\newcommand{\thema}{Vektorgeometrie \uproman{2}}
\newcommand{\lehrer}{Jonas Guler}
\newcommand{\klasse}{3GaLaWa}
\newcommand{\serie}{A}
\newcommand{\data}{\today}
\newcommand{\dauer}{45 Minuten}
\newcommand{\name}{\bf Name:}
\newcommand{\nota}{Note:}
\newcommand{\datum}{15.11.22}
%=========================================================
\begin{document}
\fontsize{14}{14}\selectfont
%========================================================

%========================================================
\begin{tabular*}{\textwidth}{l @{\extracolsep{\fill}}l @{\extracolsep{6pt}}l}
\textbf{\thema} & \textbf{Note:} & {\answerbox}\\
\textbf{Klasse \klasse} & & \\
\textbf{\lehrer}  & {\name} &\textbf{\hrulefill}\\[8pt]
%-------------------
\end{tabular*}
%--------------------
\begin{center}
\rule[1ex]{\textwidth}{1pt}
Dauer: \dauer \hspace{9cm}Datum: \datum\\
\vspace{5pt}
\rule[2ex]{\textwidth}{1pt}
\end{center}
%=========================================================

%=========================================================
\begin{center}
\textbf{Bewertungstabelle (Wird vom Lehrer ausgefüllt)}\\
\vspace{5pt}
\addpoints
\gradetable[h][questions]\\
\vspace{15pt}
Der Rechenweg muss \textbf{vollständig und nachvollziehbar} sein.\\
Lösen Sie die Prüfung mit einen schwarzen oder blauen Kugelschreiber. Resultate in Bleistift, roter oder grüner Farbe werden nicht berücksichtigt.\\
Die Schlussresultate sind auf drei Nachkommastellen zu runden oder vollständig gekürzt anzugeben.
\end{center}

\noindent
\rule[1ex]{\textwidth}{1pt}

\vspace{10cm}
\begin{center}
    %\Huge{Serie \serie}
\end{center}

\newpage

\begin{questions}
%=========================================================
%---------------------TESTFRAGEN
%=========================================================
\question[4]
\begin{parts}
\part Erkläre in eigenen Worten den Begriff \glqq Vektor\grqq.
\end{parts}
\vspace{5pt}
\antwort{8}
\begin{parts}
\setcounter{partno}{1}
\part Erkläre in eigenen Worten den Beriff \glqq kollinear\grqq\: und mache ein konkretes Zahlenbeispiel.
\end{parts}
\vspace{5pt}
\antwort{8}


%---------------------------------------------------------
\newpage
\question[4]
Berechne die Koordinaten des Schwerpunkts $S$ des Dreiecks $ABC$, wenn die folgenden Punkte gegeben sind:
\[A=(6,1,-3),\:B=(7,-7,4) \text{ und } C=(-4,0,5)\] 
\vspace{5pt}
\antwort{20}

%---------------------------------------------------------
\newpage
\question[4]
Die Länge der Strecke $\overline{AB}$ mit $A=(7,1,5)$ und $B=(6,y,-3)$ beträgt $9$. Bestimme $y$.
\vspace{5pt}
\antwort{20}

%---------------------------------------------------------
\newpage
\question[6]
Zerlege den Vektor $\vec{d}$ nach den Vektoren $\vec{a}$, $\vec{b}$ und $\vec{c}$. \[ \vec{a} = \vvector{2}{1}{3}, \vec{b} = \vvector{3}{1}{-1}, \vec{c} = \vvector{5}{-2}{1}, \vec{d} = \vvector{8}{7}{2} \]
\vspace{5pt}
\antwort{20}

%---------------------------------------------------------
\newpage
\question[3]
\begin{parts}
\part Berechne das Skalarprodukt $\Vec{a}\cdot\Vec{b}$, wenn gilt $\lvert\Vec{a}\rvert=6,\lvert\Vec{b}\rvert=4$ und $\alpha=60^\circ$.
\end{parts}
\vspace{5pt}
\antwort{4}
\begin{parts}
\setcounter{partno}{1}
\part Berechne den Zwischenwinkel von $\Vec{a}$ und $\Vec{b}$, wenn gilt $\Vec{a}\cdot\Vec{b}=10$, $\lvert\Vec{a}\rvert=2\cdot\sqrt{2}$ und $\lvert\Vec{b}\rvert=5$.
\end{parts}
\vspace{5pt}
\antwort{4}
\begin{parts}
\setcounter{partno}{2}
\part Berechne das Skalarprodukt $\Vec{a}\cdot\Vec{b}$, wenn gilt $\Vec{a}=\vvector{3}{-4}{5}$ und $\Vec{b}=\vvector{5}{2}{-1}$.
\end{parts}
\vspace{5pt}
\antwort{8}



%---------------------------------------------------------
\newpage
\question[8]
Die Vektoren $\vec{a}=\vvector{x}{4}{8}$ und $\vec{b}=\vvector{8}{y}{1}$ spannen ein Quadrat auf.\\ Bestimme $x$ und $y$.
\vspace{5pt}
\antwort{20}



\addpoints
\end{questions}
\end{document}